\begin{figure}[tbp]
\centering
\begin{tikzpicture}
\begin{groupplot}[group style={group size=1 by 2,vertical sep=1.3cm},
 width=.92\textwidth,height=6cm,xmin=8,xmax=62,xtick={10,20,30,60},
 xlabel={Window duration (min)},ymajorgrids=true,grid style={gray!20},
 legend style={at={(.5,1.03)},anchor=south,legend columns=2,column sep=6pt}]
\nextgroupplot[ylabel={Window F1 (\%)},ymin=50,ymax=100]
\addplot[ChromecastBlue,mark=*,thick] table[x=minutes,y=f1,col sep=comma]{generated/duration_UNSW-Chromecast.csv};
\addplot[SamsungOrange,mark=square*,thick,dashed] table[x=minutes,y=f1,col sep=comma]{generated/duration_UNSW-Samsung-Camera.csv};
\legend{Chromecast,Samsung camera}
\nextgroupplot[ylabel={Median alert delay (min)},ymin=0,ymax=40]
\addplot[ChromecastBlue,mark=*,thick] table[x=minutes,y=delay,col sep=comma]{generated/duration_UNSW-Chromecast.csv};
\addplot[SamsungOrange,mark=square*,thick,dashed] table[x=minutes,y=delay,col sep=comma]{generated/duration_UNSW-Samsung-Camera.csv};
\end{groupplot}
\end{tikzpicture}
\caption{Static development results on common observed support, with fresh
fitting and calibration at every duration. Window F1 uses that duration's labels.
Delay runs from attack onset to first overlapping alert-window close and is
conditional on detection. Lines connect four measured settings.}
\label{fig:duration}
\end{figure}
